\chapter{From the Equivalence Principle to Tensors}

General relativity is commonly shortened to GR, although special
relativity is rarely called SR. The abbreviation is less important than
the decision to begin at the beginning. The most direct route is the
route Einstein took: start from a physical observation simple enough to
be felt in an elevator, press it until its consequences become
unavoidable, and only then introduce the geometry.

The first observation is the equivalence principle. Gravity and
acceleration are locally indistinguishable. The word locally will
eventually carry nearly the entire burden of the statement. Before
qualifying it, however, we should see what the unqualified statement can
already do.

\section{The Elevator as a Coordinate Transformation}

Textbooks often replace Einstein's elevator by a rocket ship. An
elevator is preferable here because its acceleration is familiar. Let
\(z\) measure height in a frame attached to the ground, and let \(z'\)
measure height from the elevator floor. If the floor has height \(L(t)\),
the same event has coordinates
\begin{equation}
 z'=z-L(t),\qquad t'=t,\qquad x'=x.
\label{eq:elevator-map}
\end{equation}
The horizontal coordinate is unchanged because the elevator does not
slide sideways. Taking \(t'=t\) deliberately suppresses special
relativity for the moment.

\begin{classroomqa}
\audiencequestion{How can a course on general relativity begin by
ignoring special relativity?}
\lecturerresponse{There is a regime in which gravity matters while all
velocities remain small compared with \(c\). The gravity--acceleration
connection can therefore be isolated in Newtonian mechanics and then
combined with Minkowski geometry. This is a staging choice, not a claim
that the final theory is nonrelativistic. \lecturetimestamp{00:05:33}}
\end{classroomqa}

For two inertial frames the elevator moves uniformly,
\begin{equation}
 L(t)=vt,\qquad z'=z-vt.
\label{eq:galilean-elevator}
\end{equation}
The origins have been chosen to coincide at \(t=0\). The board keeps the
elevator sketch beside all three coordinate relations so that the
symbols never lose their physical meaning.

\begin{figure}[H]
\centering
\includegraphics[width=\linewidth]{lecture_01_figure_01.png}
\caption{The elevator coordinates after the uniform-velocity
specialization \(L(t)=vt\).}
\par\smallskip
\textit{Lecture frame: 00:09:10.}
\end{figure}

Consider only the vertical component of Newton's law,
\begin{equation}
 F=m\ddot z.
\label{eq:newton-unprimed}
\end{equation}
The force may come from a rope, an electric field, or any other ordinary
interaction. Differentiating \eqref{eq:galilean-elevator} twice gives
\(\ddot z'=\ddot z\), and hence
\begin{equation}
 F=m\ddot z'.
\label{eq:newton-primed-inertial}
\end{equation}
The form of the law is unchanged. This elementary calculation exposes
the operative content of inertial-frame equivalence: a linear
coordinate transformation takes uniform motion to uniform motion.

\begin{figure}[H]
\centering
\includegraphics[width=\linewidth]{lecture_01_figure_02.png}
\caption{Newton's equation in the two uniformly moving coordinate
frames.}
\par\smallskip
\textit{Lecture frame: 00:11:50.}
\end{figure}

Now accelerate the elevator upward at constant \(g\):
\begin{equation}
 L(t)=\frac12gt^2,\qquad
 \dot L=gt,\qquad
 \ddot L=g.
\end{equation}
The coordinate transformation becomes
\begin{equation}
 z'=z-\frac12gt^2,
\qquad
 \ddot z'=\ddot z-g.
\label{eq:accelerated-map}
\end{equation}
Substitution into Newton's law gives
\begin{equation}
 F=m(\ddot z'+g),
\qquad
 m\ddot z'=F-mg.
\label{eq:fictitious-gravity}
\end{equation}
The derivation includes a useful live sign check: the \(+mg\) first
appears on the acceleration side and becomes \(-mg\) only when moved to
the force side. The final direction is downward in the frame whose
floor accelerates upward.

\begin{figure}[H]
\centering
\includegraphics[width=\linewidth]{lecture_01_figure_03.png}
\caption{The accelerated transformation and the resulting fictitious
force \(-mg\).}
\par\smallskip
\textit{Lecture frame: 00:15:35.}
\end{figure}

The new term looks gravitational for a precise reason: it is
proportional to mass. Since inertial response is also proportional to
mass, the mass cancels from the equation of motion. A freely falling
body follows the same trajectory regardless of its mass. The orbit of
the Earth, for example, is fixed by its position and velocity in the
Sun's field, not by the Earth's own mass in the test-body
approximation.

Before Einstein this resemblance could be dismissed as a curiosity.
Einstein instead promoted it to a principle: no local mechanical experiment
can distinguish a homogeneous gravitational field from a suitably
accelerated frame.

\section{Straight Lines, Curved Coordinates, and Light}

Plot \(t\) vertically and \(z\) horizontally. For
\(z'=z-vt\), lines of constant \(z'\) are parallel straight lines.
Straight worldlines remain straight because the coordinate
transformation is linear.

\begin{figure}[H]
\centering
\includegraphics[width=\linewidth]{lecture_01_figure_04.png}
\caption{A linear spacetime grid for uniformly moving coordinates.}
\par\smallskip
\textit{Lecture frame: 00:19:45.}
\end{figure}

For the accelerated map, a line of constant \(z'\) satisfies
\begin{equation}
 z=z'+\frac12gt^2.
\end{equation}
It is a parabola lying on its side. Extending the motion backward in
time gives an elevator that descends while slowing, turns, and then
accelerates upward. Every neighboring constant-\(z'\) line is the same
parabola shifted horizontally.

\begin{figure}[H]
\centering
\includegraphics[width=\linewidth]{lecture_01_figure_05.png}
\caption{Constant-\(z'\) curves for the uniformly accelerated
coordinate system.}
\par\smallskip
\textit{Lecture frame: 00:22:25.}
\end{figure}

This is the first geometric clue. Special relativity uses linear
Lorentz transformations, which carry straight spacetime lines to
straight spacetime lines. To represent acceleration, and therefore to
imitate gravity, one needs curvilinear transformations of spacetime
coordinates.

Einstein used this apparently modest point to answer a question that
was not modest at all: does gravity affect light? At \(t=0\), let the
elevator and ground frames be instantaneously at rest relative to one
another, and send a horizontal light ray across the elevator. In the
unaccelerated frame,
\begin{equation}
 x=ct,\qquad z=0.
\end{equation}
Using \(z=z'+\tfrac12gt^2\), the same ray in elevator coordinates obeys
\begin{equation}
 z'=-\frac12gt^2
    =-\frac{g}{2c^2}x^2.
\label{eq:accelerated-light}
\end{equation}
The light ray is straight in the inertial frame and curved in the
accelerated frame. If acceleration and gravity are locally equivalent,
a gravitational field must bend light. The argument establishes the
effect before a field equation for gravity has been written.

\begin{figure}[H]
\centering
\includegraphics[width=\linewidth]{lecture_01_figure_06.png}
\caption{The horizontal light ray and its curved trajectory in the
accelerated frame.}
\par\smallskip
\textit{Lecture frame: 00:28:10.}
\end{figure}

\section{What Cannot Be Transformed Away}

An accelerated frame can manufacture an apparent uniform gravitational
field. It does not follow that every gravitational field is a
coordinate artifact. The field of a star points radially inward. A
small freely falling laboratory near one point can erase the local
acceleration, but the same transformation cannot erase the field on all
sides of the star.

\begin{figure}[H]
\centering
\includegraphics[width=\linewidth]{lecture_01_figure_07.png}
\caption{A radial field and an extended falling body used to expose
tidal acceleration.}
\par\smallskip
\textit{Lecture frame: 00:32:50.}
\end{figure}

The obstruction is tidal force. Imagine a two-thousand-mile person
falling feet first. His feet are closer to the source and accelerate
more strongly than his head, so he is stretched. Turn him sideways and
the inward directions on opposite sides have components toward one
another, so he is squeezed. In Newtonian language, neighboring points
separated by \(\xi^j\) acquire relative acceleration
\begin{equation}
 \delta a^i \simeq \partial_j g^i\,\xi^j.
\label{eq:newtonian-tide}
\end{equation}
A coordinate acceleration can be removed. Relative acceleration across
an extended freely falling body cannot.

\begin{classroomqa}
\audiencequestion{Could different rocket engines at the head and feet
simulate the stretching and therefore simulate a tidal field?}
\lecturerresponse{Applied engines can certainly stress a body, but that
is a real non-gravitational force acting differently on its parts. A
coordinate change alone cannot create invariant relative acceleration
among neighboring freely falling particles. The diagnostic concerns
unforced free fall, not a body driven by separate rockets.
\lecturetimestamp{00:36:51}}
\end{classroomqa}

Curved coordinates can nevertheless make fictitious fields look quite
complicated. A frame may accelerate vertically while oscillating
horizontally, producing a time-dependent apparent field with both
components. Even ordinary polar coordinates provide an instructive
example. A free particle moving on a straight Cartesian line first has
decreasing \(r(t)\), reaches a minimum radius, and then has increasing
\(r(t)\). Someone given only \(r(t)\) and \(\theta(t)\) might infer an
outward repulsive force. In rotating variants of these coordinates,
centrifugal and Coriolis terms appear. They are proportional to mass and
therefore look gravitational, but they leave no invariant tidal strain
on a freely falling body.

\begin{classroomqa}
\audiencequestion{What happens if the motion is described in angular or
polar coordinates?}
\lecturerresponse{The nonlinear coordinate description makes a straight
free trajectory look as though it is repelled from \(r=0\). That
effective centrifugal force is a coordinate force. Returning to an
inertial Cartesian frame removes it and reveals no tidal distortion.
\lecturetimestamp{00:37:56}}
\end{classroomqa}

\begin{figure}[H]
\centering
\includegraphics[width=\linewidth]{lecture_01_figure_08.png}
\caption{The polar-coordinate description of a straight trajectory,
with \(r(t)\) and \(\theta(t)\) marked on the board.}
\par\smallskip
\textit{Lecture frame: 00:40:20.}
\end{figure}

\begin{classroomqa}
\audiencequestion{Could several gravitating bodies cancel one another
and leave no tidal force?}
\lecturerresponse{Forces can cancel at a chosen point, but cancellation
of the field value does not generally cancel its spatial derivatives.
Near any one source the neighboring trajectories still acquire
different accelerations. A true absence of curvature requires the tidal
diagnostic to vanish throughout the region, not merely the net force at
one point. \lecturetimestamp{00:42:16}}
\end{classroomqa}

\begin{classroomqa}
\audiencequestion{How is a tidal force measured operationally?}
\lecturerresponse{Release an extended object, such as a crystal, in free
fall and look for stresses and strains that stretch it in one direction
and compress it in another. Those deformations are measurable and do
not disappear when coordinates are renamed. \lecturetimestamp{00:44:38}}
\end{classroomqa}

\begin{classroomqa}
\audiencequestion{Would a two-thousand-mile person on a merry-go-round
feel a radial tidal effect?}
\lecturerresponse{A person strapped to the carousel feels supporting
forces and may indeed get a headache. The test object must instead be
freely falling. Its trajectory looks complicated in rotating
coordinates but remains an unstressed straight trajectory in the
inertial frame, so the invariant tidal force is zero.
\lecturetimestamp{00:45:07}}
\end{classroomqa}

\begin{classroomqa}
\audiencequestion{Can the falling body be tilted so that stretching and
squeezing cancel?}
\lecturerresponse{No single orientation removes the full deformation.
The radial field stretches along one principal direction and compresses
along transverse directions. A body with finite extent in all
directions experiences the complete pattern. \lecturetimestamp{00:47:10}}
\end{classroomqa}

\begin{figure}[H]
\centering
\includegraphics[width=\linewidth]{lecture_01_figure_09.png}
\caption{Radial stretching and transverse squeezing displayed together
for an extended freely falling body.}
\par\smallskip
\textit{Lecture frame: 00:47:55.}
\end{figure}

\begin{classroomqa}
\audiencequestion{Do tidal forces mean that the equivalence principle
fails?}
\lecturerresponse{They determine how carefully it must be stated. In a
system small enough that the gravitational field varies negligibly
across it, gravity can be removed to the relevant accuracy by a freely
falling frame. The residual variation across a larger system is the
tidal field. \lecturetimestamp{00:48:08}}
\end{classroomqa}

\section{From Minkowski Geometry to a Local Metric}

We now take a temporary rest from gravity. Minkowski had already taught
that spacetime possesses an invariant local geometry. In one spatial
dimension, with \(c=1\),
\begin{equation}
 d\tau^2=dt^2-dx^2.
\label{eq:minkowski-proper-time}
\end{equation}
One may instead define \(ds^2=dx^2-dt^2=-d\tau^2\). The difference is a
signature convention, not a difference in physics.

\begin{classroomqa}
\audiencequestion{Where does the speed of light belong if it is not set
to one?}
\lecturerresponse{If \(t\) is measured in time units and \(x\) in length
units, one convenient form is
\(d\tau^2=dt^2-dx^2/c^2\). Equivalently, multiply by \(c^2\) and use
\(c^2d\tau^2=c^2dt^2-dx^2\). Choosing \(ct\) or \(x/c\) as a coordinate
absorbs the conversion factor. \lecturetimestamp{00:51:02}}
\end{classroomqa}

\begin{figure}[H]
\centering
\includegraphics[width=\linewidth]{lecture_01_figure_10.png}
\caption{The proper-time interval that supplies the geometric bridge
from special to general relativity.}
\par\smallskip
\textit{Lecture frame: 00:51:10.}
\end{figure}

Einstein recognized that the problem of deciding whether a
gravitational field is removable resembles a problem studied by Gauss
for surfaces and by Riemann in arbitrary dimensions: given a local
distance law, how can one decide whether the geometry is genuinely
curved?

In Euclidean Cartesian coordinates, Pythagoras gives
\begin{equation}
 ds^2=(dx^1)^2+(dx^2)^2+\cdots
     =\delta_{mn}\,dx^m dx^n.
\label{eq:euclidean-line-element}
\end{equation}
The index range is the dimension of the space. Nothing in the notation
changes whether there are two, three, or twenty-six coordinates.

\begin{figure}[H]
\centering
\includegraphics[width=\linewidth]{lecture_01_figure_11.png}
\caption{The Euclidean line element as a sum of squared coordinate
increments.}
\par\smallskip
\textit{Lecture frame: 00:55:00.}
\end{figure}

Place arbitrary coordinates \(x^m\) on an arbitrary surface. The
distance between neighboring points is the quadratic form
\begin{equation}
 ds^2=g_{mn}(x)\,dx^m dx^n.
\label{eq:general-line-element}
\end{equation}
The functions \(g_{mn}(x)\) are the components of the metric tensor.
The same formula also applies to a flat surface described by
curvilinear coordinates; nonconstant metric components do not by
themselves prove intrinsic curvature.

\begin{figure}[H]
\centering
\includegraphics[width=\linewidth]{lecture_01_figure_12.png}
\caption{A general metric quadratic form beside a curved coordinate
net.}
\par\smallskip
\textit{Lecture frame: 00:58:20.}
\end{figure}

For a sphere of radius \(R\), when \(\theta\) is measured from a pole,
\begin{equation}
 ds^2=R^2\left(d\theta^2+\sin^2\theta\,d\phi^2\right).
\label{eq:sphere-polar}
\end{equation}
The same longitude increment covers a large distance near the equator
and a small distance near a pole; that position dependence is exactly
what the coefficient records.

\begin{classroomqa}
\audiencequestion{Ordinary latitude is measured from the equator, so
should the coefficient be a cosine rather than a sine?}
\lecturerresponse{Yes. If \(\lambda\) denotes latitude measured from the
equator, then
\(ds^2=R^2(d\lambda^2+\cos^2\lambda\,d\phi^2)\).
Equation \eqref{eq:sphere-polar} uses the polar angle
\(\theta=\pi/2-\lambda\), measured from the north pole.
\lecturetimestamp{01:01:03}}
\end{classroomqa}

\begin{figure}[H]
\centering
\includegraphics[width=\linewidth]{lecture_01_figure_13.png}
\caption{The sphere metric and the latitude--longitude geometry that
fixes its position-dependent coefficient.}
\par\smallskip
\textit{Lecture frame: 01:00:30.}
\end{figure}

More general coordinates can introduce \(d\theta\,d\phi\) cross terms,
but the local distance remains quadratic in the differentials. Cubic
or linear terms do not belong to a Riemannian line element.

\section{Local Distance and Intrinsic Curvature}

\begin{classroomqa}
\audiencequestion{Does the metric formula also give the distance
between two widely separated points?}
\lecturerresponse{Not by direct substitution. It gives the distance
between neighboring points. A macroscopic distance requires a path and
an integral of \(ds\); if distance means the shortest path, that path
must itself be found from the geometry. Several locally shortest paths
may even connect the same endpoints. \lecturetimestamp{01:02:30}}
\end{classroomqa}

On a bumpy surface, imagine pegs at two distant points and a taut string
between them. The string may pass around one side of a hill or another,
and moving a hill changes the minimizing route. Global distance depends
on the geometry in between. Local distance is simpler: the metric
supplies a small stick at every point, and the geometry is assembled
from those infinitesimal sticks.

\begin{figure}[H]
\centering
\includegraphics[width=\linewidth]{lecture_01_figure_14.png}
\caption{The global-distance problem: a taut route must negotiate the
surface between its endpoints.}
\par\smallskip
\textit{Lecture frame: 01:04:00.}
\end{figure}

If the local pieces can be laid onto a plane without stretching or
tearing, the intrinsic geometry is flat. Pieces with the lengths of a
triangulated sphere cannot all be laid flat while preserving their
edges. The metric therefore poses a precise problem: given
\(g_{mn}(x)\), can one determine whether its infinitesimal pieces fit
together as a flat space?

\begin{figure}[H]
\centering
\includegraphics[width=\linewidth]{lecture_01_figure_15.png}
\caption{The metric and local coordinate patch interpreted as the
infinitesimal pieces from which geometry is assembled.}
\par\smallskip
\textit{Lecture frame: 01:05:20.}
\end{figure}

\begin{classroomqa}
\audiencequestion{Could one diagnose curvature by measuring the ratio
of circumference to diameter?}
\lecturerresponse{Yes, that is a useful two-dimensional test. On a
curved surface the ratio can differ from \(\pi\). For an infinitesimal
circle the deviation enters at higher order, however, and the method
does not generalize as cleanly as a curvature tensor.
\lecturetimestamp{01:07:26}}
\end{classroomqa}

A region is flat if there exists a coordinate system in which
\begin{equation}
 g_{mn}=\delta_{mn},
\qquad
 ds^2=\delta_{mn}\,dx^m dx^n
\label{eq:flat-metric}
\end{equation}
throughout that region. A complicated \(g_{mn}(x)\) may therefore
describe either a curved space or merely curved coordinates on a flat
space. The question is whether a coordinate transformation can remove
all of the complication.

\begin{figure}[H]
\centering
\includegraphics[width=\linewidth]{lecture_01_figure_16.png}
\caption{The metric problem stated on a local coordinate patch: can its
curvilinear form be transformed to the flat one?}
\par\smallskip
\textit{Lecture frame: 01:07:15.}
\end{figure}

\begin{classroomqa}
\audiencequestion{Does the spherical line element assume that each
small interval is flat even though the sphere is not?}
\lecturerresponse{To leading order, every smooth surface is approximated
at a point by its tangent plane. Curvature appears in how neighboring
tangent approximations fail to fit together, or in higher-order changes
across a finite patch. Local Euclidean behavior and global curvature
are therefore compatible. \lecturetimestamp{01:10:01}}
\end{classroomqa}

\begin{classroomqa}
\audiencequestion{Must the metric be continuous, and what happens at a
crease or crevice?}
\lecturerresponse{The discussion assumes coordinate maps and a metric
smooth enough for the required derivatives. A sheet folded without
stretching has the same intrinsic distances on both sides of the fold;
its dramatic appearance is extrinsic. A genuinely singular or
nonsmooth geometry requires separate care and is outside the smooth
setup used here. \lecturetimestamp{01:11:25}}
\end{classroomqa}

This distinction is intrinsic geometry. Bending a page in three
dimensional space changes its embedding but not distances measured
within the page. A creature confined to the page, unable to look out
into the embedding space, can still measure its intrinsic geometry.

\begin{figure}[H]
\centering
\includegraphics[width=\linewidth]{lecture_01_figure_17.png}
\caption{A folded sheet used to separate intrinsic distance from
extrinsic appearance.}
\par\smallskip
\textit{Lecture frame: 01:13:00.}
\end{figure}

\begin{classroomqa}
\audiencequestion{Why not simply diagonalize the metric?}
\lecturerresponse{At any one point a suitable coordinate choice can
diagonalize the metric and normalize it to the local flat form. The
obstruction is doing so throughout a region with one coordinate system.
Curvature cannot be detected by inspecting a single point alone.
\lecturetimestamp{01:13:48}}
\end{classroomqa}

\begin{figure}[H]
\centering
\includegraphics[width=\linewidth]{lecture_01_figure_18.png}
\caption{The flat Cartesian metric \(\delta_{mn}\) written beneath the
general local quadratic form.}
\par\smallskip
\textit{Lecture frame: 01:16:00.}
\end{figure}

The parallel with gravity is now exact rather than suggestive. A
coordinate force is removable just as a curvilinear description of
flat space is removable. Tidal acceleration is measured in general
relativity by curvature. In later notation, neighboring freely falling
worldlines satisfy the geodesic-deviation equation
\begin{equation}
 \frac{D^2\xi^\mu}{D\tau^2}
 =-R^\mu{}_{\nu\rho\sigma}
   u^\nu\xi^\rho u^\sigma.
\label{eq:geodesic-deviation-preview}
\end{equation}
Vanishing curvature means no invariant tidal field; nonzero curvature
means that no coordinate choice can remove gravity throughout the
region.

\section{Two Coordinate Systems and One Vector}

To decide whether a metric can be transformed to \(\delta_{mn}\), we
must first know how geometric objects transform. Let \(x^m\) and
\(y^m\) be two smooth coordinate systems on the same region:
\begin{equation}
 x^m=x^m(y),\qquad y^m=y^m(x).
\label{eq:mutual-coordinate-maps}
\end{equation}
The symbols \(x\) and \(y\) label complete coordinate systems, not two
particular spatial directions.

\begin{figure}[H]
\centering
\includegraphics[width=\linewidth]{lecture_01_figure_19.png}
\caption{The mutually related coordinate maps \(x^m(y)\) and
\(y^m(x)\).}
\par\smallskip
\textit{Lecture frame: 01:20:35.}
\end{figure}

For neighboring points, the chain rule gives
\begin{equation}
 dy^m=\frac{\partial y^m}{\partial x^p}\,dx^p.
\label{eq:differential-transform}
\end{equation}
The repeated index \(p\) is summed. The displacement itself is one
geometric vector; \(dx^p\) and \(dy^m\) are its components in two
coordinate systems.

\begin{figure}[H]
\centering
\includegraphics[width=\linewidth]{lecture_01_figure_20.png}
\caption{The differential coordinate transformation obtained from the
chain rule.}
\par\smallskip
\textit{Lecture frame: 01:22:30.}
\end{figure}

\begin{classroomqa}
\audiencequestion{On a curvilinear grid, what do the coordinate lines
represent?}
\lecturerresponse{In two dimensions, a line is obtained by holding one
coordinate fixed while varying the other. In higher dimensions the
corresponding constant-coordinate sets are hypersurfaces. Their local
intersections define the coordinate directions used to resolve vector
components. \lecturetimestamp{01:25:20}}
\end{classroomqa}

A contravariant vector is defined by the same transformation law as a
coordinate displacement:
\begin{equation}
 (V')^m(y)
 =\frac{\partial y^m}{\partial x^p}\,V^p(x).
\label{eq:contravariant-vector}
\end{equation}
Velocity along a surface is an example. The vector does not change when
coordinates change; its components do.

\begin{figure}[H]
\centering
\includegraphics[width=\linewidth]{lecture_01_figure_21.png}
\caption{The boxed transformation law for a contravariant vector.}
\par\smallskip
\textit{Lecture frame: 01:29:30.}
\end{figure}

Einstein's summation convention suppresses the summation sign whenever
an index is repeated once upstairs and once downstairs. The convention
made formulas easier to print, but more importantly it makes the
contraction pattern visible.

\begin{classroomqa}
\audiencequestion{Are special assumptions being made about which
coordinate systems are allowed?}
\lecturerresponse{Only that the two systems are functionally related
and smooth enough for the displayed derivatives and inverse map to
exist in the region under discussion. No linearity or orthogonality is
assumed. \lecturetimestamp{01:28:12}}
\end{classroomqa}

\section{Covectors and the Chain Rule}

Let \(S(x)\) be a scalar: it has one well-defined value at each point,
independent of the coordinate labels. Temperature and a scalar field
are examples. Its gradient has components
\begin{equation}
 W_p=\frac{\partial S}{\partial x^p}.
\label{eq:gradient-covector}
\end{equation}

\begin{figure}[H]
\centering
\includegraphics[width=\linewidth]{lecture_01_figure_22.png}
\caption{A scalar gradient introduced beside the contravariant
transformation law.}
\par\smallskip
\textit{Lecture frame: 01:31:40.}
\end{figure}

Applying the chain rule in the \(y\) coordinates,
\begin{equation}
 W'_m
 =\frac{\partial S}{\partial y^m}
 =\frac{\partial x^p}{\partial y^m}
  \frac{\partial S}{\partial x^p}
 =\frac{\partial x^p}{\partial y^m}\,W_p.
\label{eq:covariant-vector}
\end{equation}
This is the transformation law of a covariant vector, or covector. The
Jacobian is the inverse of the one in
\eqref{eq:contravariant-vector}; lower and upper indices therefore
encode different transformation behavior.

\begin{figure}[H]
\centering
\includegraphics[width=\linewidth]{lecture_01_figure_23.png}
\caption{Contravariant and covariant transformation laws displayed
together.}
\par\smallskip
\textit{Lecture frame: 01:34:40.}
\end{figure}

\begin{classroomqa}
\audiencequestion{What is the connection between covariant and
contravariant vectors?}
\lecturerresponse{At this stage they are defined as different kinds of
objects: displacements and velocities transform one way, gradients the
inverse way. The metric will later provide a one-to-one map that lowers
or raises an index. That relationship should not be assumed before the
metric is introduced. \lecturetimestamp{01:37:43}}
\end{classroomqa}

\begin{classroomqa}
\audiencequestion{Did lowering an index previously change the sign of
the time component?}
\lecturerresponse{That sign came from the Lorentzian metric. The present
derivation is first being done in ordinary Euclidean geometry. Raising
or lowering always uses the relevant metric; a time-component sign
appears only when that metric has Lorentzian signature.
\lecturetimestamp{01:38:29}}
\end{classroomqa}

\begin{classroomqa}
\audiencequestion{Should the two transformation laws already be viewed
as two component descriptions of the same vector?}
\lecturerresponse{Not yet. Once the metric is available, a
contravariant vector \(V^m\) can be converted to a covector
\(V_m=g_{mn}V^n\). Until that map has been established, upper- and
lower-index objects are logically distinct. \lecturetimestamp{01:39:00}}
\end{classroomqa}

These formulas are independent of dimension. The index range changes,
but the transformation laws do not.

\section{Tensors with More Than One Index}

Take two contravariant vectors \(V^m\) and \(U^n\). Their product
\begin{equation}
 T^{mn}=V^mU^n
\label{eq:rank-two-product}
\end{equation}
is a rank-two contravariant tensor. In three dimensions it has nine
components; in four dimensions it has sixteen. The order of the indices
matters unless a separate symmetry is given: \(T^{mn}\) need not equal
\(T^{nm}\).

\begin{figure}[H]
\centering
\includegraphics[width=\linewidth]{lecture_01_figure_24.png}
\caption{A rank-two tensor built as the product of two contravariant
vectors.}
\par\smallskip
\textit{Lecture frame: 01:43:40.}
\end{figure}

Each upper index contributes one forward Jacobian:
\begin{equation}
 (T')^{mn}
 =\frac{\partial y^m}{\partial x^p}
  \frac{\partial y^n}{\partial x^q}
  T^{pq}.
\label{eq:rank-two-contravariant}
\end{equation}
The product construction motivates the rule, but the transformation
law is the definition of a tensor. A rank-three contravariant tensor
would acquire three such Jacobian factors.

\begin{figure}[H]
\centering
\includegraphics[width=\linewidth]{lecture_01_figure_25.png}
\caption{The two Jacobian factors in the transformation of
\(T^{mn}\).}
\par\smallskip
\textit{Lecture frame: 01:44:20.}
\end{figure}

Likewise, each lower index contributes an inverse Jacobian:
\begin{equation}
 T'_{mn}
 =\frac{\partial x^p}{\partial y^m}
  \frac{\partial x^q}{\partial y^n}
  T_{pq}.
\label{eq:rank-two-covariant}
\end{equation}
A mixed tensor uses one factor of each kind.

\begin{figure}[H]
\centering
\includegraphics[width=\linewidth]{lecture_01_figure_26.png}
\caption{The completed transformation law for a tensor with two
covariant indices.}
\par\smallskip
\textit{Lecture frame: 01:47:30.}
\end{figure}

The lecture ends just as this machinery returns to the metric. Since
the interval is a scalar, it must have the same value in either
coordinate system:
\begin{align}
 ds^2
 &=g_{mn}(x)\,dx^m dx^n,\\
 dx^m
 &=\frac{\partial x^m}{\partial y^p}\,dy^p.
\end{align}
Substitution completes the line that the board leaves unfinished:
\begin{equation}
 g'_{pq}(y)
 =\frac{\partial x^m}{\partial y^p}
  \frac{\partial x^n}{\partial y^q}
  g_{mn}(x(y)).
\label{eq:l1-metric-transformation}
\end{equation}
Thus \(g_{mn}\) is a rank-two covariant tensor.

\begin{figure}[H]
\centering
\includegraphics[width=\linewidth]{lecture_01_figure_27.png}
\caption{The metric transformation begun from invariance of the line
element at the close of the lecture.}
\par\smallskip
\textit{Lecture frame: 01:48:40.}
\end{figure}

The opening physical question has now become a mathematical one. Given
\(g_{mn}(x)\), does there exist a coordinate transformation for which
\(g'_{mn}=\delta_{mn}\) throughout the region? If so, the apparent
complexity is only coordinate complexity. If not, the obstruction is
curvature, and in spacetime that obstruction is the invariant content
of gravity.
